Die Tabelle (FoTaBe S.~190) gibt für Backstein die Wärmeleitfähigkeit $\lambda = \lBO$ und für eine beidseits verputzte Backsteinmauer der Dicke $d = \dMO$ den Wärmedurchgangskoeffizienten $U = \UMO$.

Die beiden Grössen haben verschiedene Einheiten und lassen sich nicht direkt vergleichen. Die Wärmeleitfähigkeit ist eine Materialeigenschaft: Der Wärmestrom durch eine Schicht ist $P = \lambda\,A\,\Delta T/d$, er hängt also noch von der Dicke ab. Der Wärmedurchgangskoeffizient beschreibt ein ganzes Bauteil, $P = U\,A\,\Delta T$, wobei $\Delta T$ der Unterschied der \emph{Lufttemperaturen} innen und aussen ist.

Für eine reine Backsteinschicht der gleichen Dicke wäre
\begin{align}
 \frac{\lambda}{d} &= \UbF = \frac{\lB}{\dM} = \Ub \approx \result{\UbP{0}{2}}
\end{align}
also mehr als der Tabellenwert der Mauer. Die Mauer lässt weniger Wärme durch, weil zur Backsteinschicht weitere Widerstände in Serie geschaltet sind: der Verputz und vor allem die Wärmeübergänge zwischen Luft und Wandoberfläche. An beiden Oberflächen bildet sich eine dünne Schicht fast ruhender Luft, und Luft ist ein sehr schlechter Wärmeleiter. Die Wärmewiderstände (pro Fläche) addieren sich:
\begin{align}
 \frac{1}{U} &= \frac{1}{h_\text{i}} + \frac{d}{\lambda} + \frac{1}{h_\text{a}} + \ldots
\end{align}
Für Verputz und Übergänge bleibt
\begin{align}
 \frac{1}{U}-\frac{d}{\lambda} &= \frac{1}{\UM}-\frac{\dM}{\lB} = \Rue \approx \result{\RueP{0}{2}}
\end{align}
Umgekehrt ist der Temperaturunterschied zwischen den beiden Oberflächen des Backsteins kleiner als der zwischen der Luft innen und aussen.
